\documentclass[aspectratio=169]{beamer}

\usetheme{Madrid}
\usecolortheme{default}

\title{Fault Injection Attack on Kyber768}
\subtitle{Clock Glitching and Instruction Skipping}
\author{Pranjal Jat}
\date{\today}

\begin{document}

\begin{frame}
    \titlepage
\end{frame}

\begin{frame}{Outline}
    \tableofcontents
\end{frame}

\section{Introduction to Fault Injection}
\begin{frame}{What is Fault Injection?}
    \begin{itemize}
        \item \textbf{Defintion}: Intentionally introducing errors into a system to alter its intended behavior.
        \item \textbf{Purpose}: 
            \begin{itemize}
                \item Evaluate system robustness (Safety).
                \item Bypass security mechanisms (Security).
            \end{itemize}
        \item \textbf{Classification of Attacks}:
            \begin{enumerate}
                \item \textbf{Invasive}: Direct physical modification (e.g., FIB, probing silicon). Expensive, requires decapsulation.
                \item \textbf{Semi-Invasive}: Access to chip surface but no depackaging (e.g., Optical Fault Injection/Lasers).
                \item \textbf{Non-Invasive}: Manipulating accessible environmental parameters (e.g., \textbf{Clock}, \textbf{Voltage}, EM, Temperature). Low cost, scalable.
            \end{enumerate}
    \end{itemize}
\end{frame}

\begin{frame}{Mechanics of Glitching}
    \begin{block}{The Core Principle: Timing Violations}
    Digital circuits rely on strict timing constraints:
    \begin{itemize}
        \item \textbf{Setup Time}: Data must be stable \textit{before} the clock edge.
        \item \textbf{Hold Time}: Data must remain stable \textit{after} the clock edge.
        \item \textbf{Propagation Delay}: Time taken for signals to travel through logic gates.
    \end{itemize}
    \textit{Glitching aims to violate these constraints, causing the CPU to latch incorrect values or execute instructions incorrectly.}
    \end{block}
\end{frame}

\begin{frame}{Clock vs. Voltage Glitching}
    \begin{columns}
        \column{0.5\textwidth}
        \textbf{Clock Glitching}
        \begin{itemize}
            \item \textbf{Method}: Temporarily shortening the clock period (higher frequency).
            \item \textbf{Effect}: The critical path logic cannot finish computation before the next clock edge arrives.
            \item \textbf{Result}: Setup time violation $\rightarrow$ Latching intermediate/wrong state.
            \item \textit{Precise, often single-cycle resolution.}
        \end{itemize}

        \column{0.5\textwidth}
        \textbf{Voltage Glitching}
        \begin{itemize}
            \item \textbf{Method}: Rapidly dropping the supply voltage (VCC).
            \item \textbf{Effect}: Lower voltage increases the \textbf{propagation delay} of transistors.
            \item \textbf{Result}: Logic is too slow for the standard clock speed $\rightarrow$ Setup time violation.
            \item \textit{Global effect, affects entire chip path.}
        \end{itemize}
    \end{columns}
\end{frame}

\begin{frame}{Common Fault Models}
    \begin{itemize}
        \item \textbf{Instruction Skip}: preventing an instruction from executing (effectively turning it into a \texttt{NOP}).
        \item \textbf{Instruction Corruption}: transforming one instruction into another (e.g., \texttt{ADD} becomes \texttt{SUB}).
        \item \textbf{Register/Memory Corruption}: flipping bits in data storage.
        \item \textbf{Control Flow Hijacking}: manipulating branch tests to take the wrong path.
    \end{itemize}
\end{frame}

\section{Kyber: Post-Quantum Cryptography}
\begin{frame}{Kyber Overview}
    \begin{itemize}
        \item \textbf{Kyber}: A lattice-based Key Encapsulation Mechanism (KEM).
        \item Selected by NIST as the primary Post-Quantum Cryptography (PQC) standard (ML-KEM).
        \item \textbf{Security}: Based on the hardness of the Module-LWE (Learning With Errors) problem.
        \item \textbf{Decapsulation Process}:
            \begin{enumerate}
                \item Decrypt ciphertext to get message $m'$.
                \item Re-encrypt $m'$ to get $c'$.
                \item \textbf{Implicit Rejection}: Check if received ciphertext $c == c'$.
                \item If equal: Return $K = \text{KDF}(m', \dots)$.
                \item If not equal: Return $K = \text{KDF}(\sigma, \dots)$ (Random).
            \end{enumerate}
    \end{itemize}
\end{frame}

\section{Attack Methodology}
\begin{frame}{The Target: Instruction Skip}
    \begin{itemize}
        \item \textbf{Objective}: Force the device to skip the check (or the result of the check) that validates the ciphertext.
        \item \textbf{Target Instruction}: \texttt{negs} (ARM assembly: conditional negation).
        \item \textbf{Mechanism}:
            \begin{itemize}
                \item In constant-time implementations (like pqm4), branching is avoided.
                \item Conditional moves (\texttt{cmov}) or arithmetic operations (\texttt{negs}) are used to select the final key.
                \item If we skip \texttt{negs} or effect the \texttt{cmov}, we can force the device to output the "True" key even when the ciphertext is invalid.
            \end{itemize}
    \end{itemize}
\end{frame}

\begin{frame}{Clock Glitching Strategy}
    \begin{itemize}
        \item \textbf{Setup}: ChipWhisperer Lite + STM32F3 (running pqm4).
        \item \textbf{Approach}:
            \begin{enumerate}
                \item \textbf{Profiling}: Used CPA (Correlation Power Analysis) to identify the exact cycle of the comparison operation (approx. cycle 22 relative to trigger).
                \item \textbf{Glitching}: Drive the target clock with \texttt{clock\_xor} glitch.
                \item \textbf{Parameters}:
                    \begin{itemize}
                        \item \textbf{Ext Offset}: Cycle count to wait before glitching.
                        \item \textbf{Width}: Duration of the glitch pulse (4.0\% - 8.0\%).
                        \item \textbf{Offset}: Phase shift to target specific edge.
                    \end{itemize}
            \end{enumerate}
    \end{itemize}
\end{frame}

\section{Implementation Details}
\begin{frame}{Implementation Logic}
    \begin{columns}
        \column{0.5\textwidth}
        \textbf{Firmware (\texttt{test\_glitch.c})}:
        \begin{itemize}
            \item Generate Keypair ($pk, sk$).
            \item Create valid $ct$.
            \item \textbf{Corrupt} $ct$ (flip last byte).
            \item Loop: Decapsulate($sk, ct_{corrupt}$).
            \item Normal Result: Random Key.
            \item \textbf{Success}: Gold Key (Check Bypassed).
        \end{itemize}

        \column{0.5\textwidth}
        \textbf{Attack Script (\texttt{attack\_negs.py})}:
        \begin{itemize}
            \item Reset Target.
            \item Wait for "WAIT" signal.
            \item Insert Glitch at parameterized time.
            \item Read Output Key.
            \item Compare with Gold Key.
        \end{itemize}
    \end{columns}
\end{frame}

\section{Challenges \& Current State}
\begin{frame}{Challenges Faced}
    \begin{itemize}
        \item \textbf{Timing Precision}: Finding the Single Instruction among millions of cycles.
            \begin{itemize}
                \item \textit{Solution}: 24 MHz CPA analysis narrowed search window to $\sim$cycle 22.
            \end{itemize}
        \item \textbf{State Corruption ("Different" Keys)}:
            \begin{itemize}
                \item Glitches often corrupt the calculation rather than skipping the instruction cleanly.
                \item Result: 1000s of keys that are neither "Gold" nor "Random" -> Just noise.
            \end{itemize}
        \item \textbf{Stability}:
            \begin{itemize}
                \item Target device hangs/crashes frequently.
            \end{itemize}
        \item \textbf{Randomness}:
            \begin{itemize}
                \item Need deterministic PRNG (`randombytes.c`) to verify results against a known "Gold" secret.
            \end{itemize}
    \end{itemize}
\end{frame}

\begin{frame}{Current State}
    \begin{itemize}
        \item \textbf{Status}: Active Refinement.
        \item Successfully identified the sensitive timing window (Cycle 20-25).
        \item \textbf{Observation}: 
            \begin{itemize}
                \item Key matches Gold Key on upper bits, but LSBs are corrupted.
                \item Indicates we are hitting the right operation but timing/width needs fine-tuning.
            \end{itemize}
        \item \textbf{Next Steps}:
            \begin{itemize}
                \item Finer sweep of Phase Offset.
            \end{itemize}
    \end{itemize}
\end{frame}

\begin{frame}
    \centering
    \Huge Thank You
\end{frame}

\end{document}
